\chapter{Strukturpolynome und die Horizontableitung}
\label{ch:strukturpolynome}

% Register: D3, D4, D7, D10, S2, S3, S4, S7, S39, S40, S41, S42, S43.
% Quellen: Spez. 2, 3, 4, 9, Anhänge H und I.

\section{Die Gestalt bekommt Koordinaten}

Bisher war die „Gestalt“ einer Hori-Zahl nur ihr Ziffernbild.
Jetzt geben wir ihr eine algebraische Form -- und damit wird aus
einer Anschauung ein Rechenwerkzeug.

Der Schlüssel sind die \emph{fallenden Fakultäten}
\[
\ff X0=1,\qquad \ff Xd=X(X-1)\cdots(X-d+1),
\]
die Polynome, die „geordnetes Auswählen ohne Zurücklegen“ zählen.
Liest man die Ziffern einer Hori-Zahl von rechts nach links als
Koeffizienten dieser Basis, entsteht ihr \emph{Strukturpolynom}:
\[
a=(a_1,\ldots,a_n)
\quad\longmapsto\quad
P_a(X)=\sum_{d=0}^{n-1}a_{n-d}\,\ff Xd .
\]

\begin{satz}[Auswertungssatz; Register S2]
$V_n(a)=P_a(n+1)$.
\end{satz}

\begin{proof}
Für $d=n-i$ ist $\ff{(n+1)}d=(n+1)!/(i+1)!$ -- genau das Gewicht
der Stelle $i$. Einsetzen liefert die Wertformel.
\end{proof}

Das ist eine kleine Rechnung mit einer großen Konsequenz: Der Wert
einer Hori-Zahl ist \emph{die Auswertung ihrer Gestalt an der Stelle
ihres Horizonts}. Die Sieben in $H_4$, also $0012$, hat die Gestalt
$P(X)=X+2$; ihr Wert ist $P(5)=7$. Stellt man eine Null voran,
bleibt die Gestalt $X+2$ -- nur der Auswertungspunkt wandert:
$P(6)=8$. Die mysteriöse Wertverschiebung durch führende Nullen ist
kein Rätsel mehr, sondern schlicht: \emph{dasselbe Polynom, eine
Stelle weiter ausgewertet.}

Zwei Vorsichtsmaßnahmen. Erstens: Zulässig sind nur nichtnegative
ganzzahlige Koeffizienten, und im Horizont $n$ nur Polynome mit
$d+c_d\le n$ an allen besetzten Stellen ($c_d>0$) --
das ist die Ziffernschranke in Polynomsprache; die Kennzahl
$\sigma(P)=\max_{c_d>0}(d+c_d)$ -- das Maximum läuft über die
besetzten Stellen, und für das Nullpolynom vereinbaren wir
$\sigma(0)=0$ -- heißt Strukturhorizont (Register D4).
Zweitens: Das Polynom allein ist \emph{nicht} die ganze Hori-Zahl.
$12$, $012$ und $0012$ haben dieselbe Gestalt $X+2$, sind aber drei
verschiedene Elemente -- erst das Paar aus Horizont und Polynom
trägt alles.

\section{Die Horizontableitung}

Wie stark verändert eine führende Null den Wert? Die Antwort ist
die diskrete Ableitung, in der Analysis als Vorwärtsdifferenz
bekannt:
\[
\Delta P(X)=P(X+1)-P(X),
\qquad
\Delta\,\ff Xd=d\,\ff X{d-1}.
\]
Die zweite Formel ist das diskrete Ebenbild von
$(x^d)'=d\,x^{d-1}$ -- die fallenden Fakultäten sind für das
Differenzenrechnen, was die Potenzen fürs Differenzieren sind.

\begin{satz}[Register S7]
Die Wertänderung bei strukturtreuer Erweiterung ist
$\Delta P(n+1)$. Insbesondere ist eine führende Null genau dann
wertneutral, wenn $\Delta P=0$, also $P$ konstant ist -- das sind
exakt die Zahlen in oder über ihrem Eigenhorizont.
\end{satz}

\begin{beispiel}
Die Gestalt $X+2$ (Ziffern $\ldots12$) hat Ableitung $1$: Ihre Werte
klettern beim Umziehen um genau eins, $5,6,7,8,\ldots$ Die Gestalt
$\ff X2$ -- das ist die Zwölf in $H_3$, Ziffern $100$ -- hat
Ableitung $2X$: Ihre Werte springen $12, 20, 30, 42,\ldots$, mit
Zuwächsen $8, 10, 12$, die selbst wachsen. Eine Konstante wie die
Gestalt $7$ hat Ableitung $0$ und rührt sich nie.
\end{beispiel}

Damit ist der Eigenhorizont aus Kapitel~\ref{ch:erweiterungen} kein
anschaulicher Begriff mehr, sondern eine Gleichung: Die
\emph{Gestalten} mit Horizontableitung null sind genau die
konstanten Polynome -- als Gestalten gelesen also die natürlichen
Zahlen. (Genau genommen: Als Elemente bilden die Konstanten viele
horizontale Vertreter desselben Werts; erst im gestalttreuen
Grenzraum aus Kapitel~\ref{ch:erweiterungen} wird daraus exakt eine
Kopie von $\Nn$.) Alles andere im Horiraum hat eine
echte „Geschwindigkeit“: Eine Gestalt vom Grad 1 wächst pro
Horizontschritt um einen konstanten Betrag, Grad 2 beschleunigt,
und so fort. Der Polynomgrad ist die \emph{Hori-Dimension} einer
Gestalt (Register D7).

\section{Strecke mal Strecke ist Fläche -- jetzt exakt}

Die Dimension verhält sich unter Multiplikation additiv:
$\dim(PQ)=\dim P+\dim Q$ für $P,Q\ne0$. Was in der Schule eine Eselsbrücke ist --
Länge mal Länge gibt Fläche --, ist hier ein Satz über Grade. Und
man kann direkt mit Gestalten rechnen. Die Vier hat in $H_3$ die
Gestalt $X$ (Ziffern $010$), die Fünf die Gestalt $X+1$ (Ziffern
$011$). Ihr Gestaltprodukt ist
\[
X(X+1)=\ff X2+2X
\]
(denn $X^2=\ff X2+X$), also die Ziffernfolge $120$ -- und
tatsächlich: $120$ in $H_3$ hat den Wert $1\cdot12+2\cdot4+0=20$.
Die Multiplikation $4\cdot5=20$ ist hier ohne einen einzigen
Übertrag geschehen, rein durch Kombinieren der Gestalten. Die
Produktformel der fallenden Fakultäten hat dabei eine hübsche
Nebenbedeutung: $X\cdot X=\ff X2+X$ zerlegt alle geordneten Paare
in die mit verschiedenen und die mit gleichen Einträgen.

\section{Was Operationen kosten: das Horizontkalkül}
\label{sec:horizontkalkuel}

Jede Gestalt braucht einen Horizont, der groß genug ist -- die
Kennzahl dafür ist der Strukturhorizont $\sigma$. Damit bekommt jede
Operation einen \emph{Preis}: Wie viel Kapazität muss der Zielraum
haben, damit das Ergebnis dort sicher Platz findet? Für einen
Operator $T$ fragen wir nach dem schlimmsten Fall
\[
\Gamma_T(r,s)=\max\{\sigma(T(P,Q)) \mid \sigma(P)\le r,\ \sigma(Q)\le s\}
\]
(Register D10; für einstellige Operatoren sinngemäß $\Gamma_T(r)$). Auf den ersten Blick ein Ungetüm: Schon für $r=s=8$
wären über hundert Milliarden Paare zu prüfen. Es geht mit einem
einzigen.

\begin{satz}[Extremalprinzip; Register S39]
In jedem Horizont $r$ gibt es die \emph{maximal gefüllte} Gestalt
$M_r$ -- jede Ziffer auf Anschlag,
$M_r(X)=\sum_{d<r}(r{-}d)\,\ff Xd$; wir nennen sie auch das
\emph{Maximalpolynom} des Horizonts.
Für jeden Operator, der auf wachsende Koeffizienten mit wachsenden
Koeffizienten antwortet (das tun alle aus Addition, Multiplikation
und $\Delta$ zusammengesetzten), gilt
\[
\Gamma_T(r,s)=\sigma\bigl(T(M_r,M_s)\bigr).
\]
\end{satz}

Der schlimmste Fall unter Abermilliarden Kandidaten ist immer
derselbe: alle Ziffern voll. Der Beweis (Anhang~\ref{ch:anhang-beweise})
ist eine Zeile Monotonie -- aber die Konsequenz ist ein
\emph{Kalkül}: Addition kostet exakt $r+s$; die Multiplikation kostet
$\beta(r,s)$, die im Kern fakultätisch wachsende Produkthöhe, die uns
in Kapitel~\ref{ch:zaehlungen} als Zählfolge wiederbegegnet. Und die
Ableitung liefert die größte Überraschung des Kalküls:

\begin{satz}[Differenzüberlauf; Register S40]
Für $r\ge2$ gilt
\[
\max_{\sigma(P)\le r}\sigma(\Delta P)
=\Bigl\lfloor\tfrac{(r+1)^2}{4}\Bigr\rfloor-1 .
\]
\end{satz}

Man halte kurz inne: $\Delta$ \emph{senkt} den Grad -- und lässt die
Kapazität trotzdem quadratisch anschwellen. Schon bei $r=8$ hebt eine
einzige Differenz den Strukturhorizont von 8 auf 19. Der Zeuge ist
eine Gestalt mit einer einzigen Ziffer in der Mitte: Die Ableitung
multipliziert jeden Koeffizienten mit seinem Grad, und am teuersten
ist der Kompromiss aus „hoher Grad“ und „großer Koeffizient“ genau
in der Mitte -- daher der Term $\lfloor(r{+}1)^2/4\rfloor$, das
größte Produkt zweier Zahlen mit fester Summe. Die Lehre: Der Grad
misst die \emph{Dimension} einer Gestalt, $\sigma$ misst ihre
\emph{Kapazität}, und die beiden können gegenläufig springen. Eine
gewöhnliche Gradbetrachtung ist für diesen Effekt blind.

Eine Operation fehlt dem Kalkül noch, die heikelste: das
\emph{Einsetzen} einer Gestalt in eine andere, $P\circ Q$. Dass das
überhaupt erlaubt ist, ist selbst ein Satz -- es ist von vornherein
nicht klar, dass beim Einsetzen wieder lauter zulässige, ganzzahlige,
nichtnegative Ziffern entstehen.

\begin{satz}[Kompositionsabschluss; Register S41]
Setzt man eine Gestalt in eine Gestalt ein, entsteht wieder eine
Gestalt. Auch das Einsetzen gehorcht dem Extremalprinzip:
$\Gamma_\circ(r,s)=\sigma(M_r\circ M_s)$.
\end{satz}

Der Beweis (Anhang~\ref{ch:anhang-beweise}) ist das kombinatorische
Herzstück dieses Kapitels: Man liest $Q$ als Bauplan-Vorrat --
jede Ziffer $a_k$ stellt $a_k$ Schablonen bereit, die auf $k$
verschiedene Plätze montiert werden --, dann zählt $P\circ Q$
Tupel solcher Montagen, und die nötige Teilbarkeit durch $m!$
kommt daher, dass das Umbenennen der Plätze keine Montage-Auswahl
festhalten kann. Die Kosten des Einsetzens sind dann allerdings von
ganz anderem Kaliber als alles Bisherige: Die Diagonale
$\Gamma_\circ(r,r)$ beginnt mit
\[
1,\quad 4,\quad 23,\quad 6541,\quad 477\,149\,751,\quad\ldots
\]
-- Einsetzen multipliziert die Grade \emph{und} lässt die
Koeffizienten explodieren; das ist die neunte und mit Abstand am
schnellsten wachsende Zählfolge dieses Projekts
(Kapitel~\ref{ch:zaehlungen}). Der harmloseste Spezialfall, das
bloße Verschieben $P(X)\mapsto P(X{+}1)$, kostet exakt
$r+\lfloor r^2/4\rfloor$ -- auf den Punkt so viel wie eine Ableitung
einen Horizont höher (Register S42), eine kleine Dualität, die das
Kalkül gratis abwirft.

\section{Was die Gestalten zählen: die Zählbreite}
\label{sec:zaehlbreite}

Zum Abschluss des Kapitels die Frage, die man einem algebraischen
Objekt immer stellen sollte: Zählt es etwas? Die Antwort kam als
externe Einordnung ins Projekt und ist die bisher beste Brücke der
Horimetrik in ein etabliertes Gebiet -- die \emph{Zählkomplexität}.

Der Schlüssel ist die Grundbedeutung der fallenden Fakultät:
$n^{\underline d}$ zählt die geordneten $d$-Tupel
\emph{verschiedener} Objekte aus einem Vorrat von $n$. Ist nun $n$
die Anzahl der Lösungen irgendeines Suchproblems -- in der
Informatik sagt man: der \emph{Zeugen} --, dann liest sich eine
Gestalt $P=\sum_d c_d\ff Xd$ als Bauplan für neue Zählobjekte:
$c_0$ feste Sonderobjekte, $c_1$ Farben für jede einzelne Lösung,
$c_2$ Varianten für jedes geordnete Paar verschiedener Lösungen,
und so fort. $P(n)$ ist exakt die Anzahl dieser gefärbten
Zeugen-Tupel. Diese harmlose Beobachtung hat eine präzise
Konsequenz:

\begin{satz}[Hori-\#P-Satz; Register S43]
Ist $f$ die Zeugen-Zählfunktion eines beliebigen NP-Problems und
$P$ eine Gestalt, so ist auch $P\circ f$ die Zeugen-Zählfunktion
eines NP-Problems. Jede Gestalt ist also ein
\emph{Abschlussoperator} der Zählklasse \#P.
\end{satz}

Der Beweis (Anhang~\ref{ch:anhang-beweise}) ist genau die
Bauplan-Lektüre: rate den Grad, rate die Farbe, rate ein Tupel
verschiedener Zeugen. Ehrlichkeit zuerst: Die \emph{qualitative}
Seite dieser Aussage ist bekannt -- Polynome mit nichtnegativen
ganzzahligen Koeffizienten in der \emph{Binomial}basis heißen in
der Literatur \emph{binomial-good} und sind als die
(relativierenden) \#P-Abschlussoperationen charakterisiert. Wegen
$\ff Xd=d!\binom Xd$ sind unsere Gestalten eine \emph{echte
Teilklasse} davon: Die Teilbarkeit durch $d!$ ist genau der
Unterschied zwischen geordneten Tupeln und ungeordneten Auswahlen.
Was die Literatur offenbar \emph{nicht} hat, ist die quantitative
Schicht -- und die liefert die Horimetrik gratis mit:

Plötzlich bedeuten alle Größen dieses Kapitels etwas. Die
\emph{Addition} von Gestalten ist die disjunkte Vereinigung zweier
Zählkonstruktionen. Die \emph{Multiplikation} koppelt zwei
Konstruktionen, und ihre Produktformel-Koeffizienten zählen exakt
die Überlappungsmuster zweier Zeugengruppen -- das
$\beta$-Wachstum misst also, wie die Überlappungsvielfalt beim
Koppeln explodiert. Die \emph{Horizontableitung} zählt, welche
neuen Strukturen ein einziger zusätzlicher Zeuge ermöglicht
($d$ Positionen mal Rest-Tupel: genau $d\,\ff n{d-1}$) -- und der
Differenzüberlauf besagt: Ein neuer Zeuge senkt zwar die
Tupelordnung, kann die Variantenvielfalt aber quadratisch aufblähen.
Und $\sigma$ selbst? Zeichnet man das Koeffizientenbrett -- Spalte
$d$ mit $c_d$ Kästchen --, so ist $\sigma(P)=\max(d+c_d)$ die
kleinste \emph{Treppe}, in die das Brett passt: das Maximum aus
„wie viele Zeugen gleichzeitig“ und „wie viele Varianten davon“.
Wir nennen diese Größe die \emph{Zählbreite} einer Konstruktion.
Jeder Horizont $r$ ist damit eine endliche, vollständig
aufzählbare Familie von genau $(r{+}1)!$
\#P-Abschlussoperatoren, und das Horizontkalkül aus
Abschnitt~\ref{sec:horizontkalkuel} beziffert optimal, wie die
Zählbreite beim Vereinigen, Koppeln, Erweitern und Einsetzen
wächst.

Zwei Sätze zur Abgrenzung, damit kein falscher Zauber entsteht: In
dieser Nachbarschaft wohnt eine gewichtige offene Vermutung (schon
ihre einfachste Fassung würde P$\,\ne\,$NP implizieren), und
unsere Operatoren liegen ausdrücklich auf der \emph{gut
verstandenen} Seite dieser Landschaft -- die Horimetrik löst hier
gar nichts. Was sie beiträgt, ist bescheidener und vielleicht
gerade deshalb brauchbar: aus einem Erlaubt-oder-nicht-erlaubt
erstmals ein \emph{Wie breit} zu machen.

\begin{oberflaeche}
Was dieses Kapitel gebracht hat, in einem Satz: Die Gestalt einer
Hori-Zahl ist ein Polynom, ihr Wert ist dieses Polynom am Ort ihres
Horizonts ausgewertet, und die rätselhafte Wirkung führender Nullen
ist einfach die Ableitung dieses Polynoms. Die gewöhnlichen Zahlen
sind die Gestalten mit Ableitung null -- die einzigen Bewohner des
Horiraums, denen es egal ist, wo sie wohnen. Alle anderen bewegen
sich, wenn ihre Welt wächst. Jede Rechenoperation hat einen exakt
bezifferbaren Preis in Kapazität, den stets dasselbe extreme
Objekt bezahlt: die randvolle Gestalt. Und zuletzt hat sich
gezeigt, dass die Gestalten etwas \emph{zählen} -- gefärbte
Gruppen verschiedener Lösungen eines Suchproblems -- und $\sigma$
die Breite dieser Zählkonstruktionen misst. Im nächsten Kapitel
lassen wir diese bewegten Objekte rechnen -- und stoßen auf die
Frage, was dabei mit ihrer Vergangenheit geschehen soll.
\end{oberflaeche}
